\documentclass[12pt]{article}

\usepackage{amsmath}
\usepackage{graphicx}
\usepackage[
  backend=biber,
  style=authoryear
]{biblatex}
\addbibresource{refs.bib}

\title{304 SS RIS and Oxidation Prediction by Off-Lattice KMC Model}
\author{Zhuo Wang}
\date{\today}
% \setlength{\parindent}{0pt}
% \setlength{\parskip}{0.8em}
\begin{document}

\maketitle

\section{Introduction}

% Diffusion flux follows:
% \begin{equation}
%     J = -D \nabla C
% \end{equation}
% Einstein derived relativity \parencite{einstein1905}.

\section{Simulation Method} \label{sec:method}

In this study, an off-lattice kinetic Monte Carlo (olKMC) model is employed to simulate radiation-induced segregation (RIS) and grain boundary oxidation in 304 stainless steel (304 SS). The simulation framework consists of two main parts: elemental diffusion and grain boundary oxidation. In the following subsections, the modeling approach for RIS and the treatment of oxidation processes are described separately. 

\subsection{Radiation Induced Segregation} \label{sec:ris}

To phenomenologically capture RIS, we empoly an effective chemical potential that accounts for the preferential enrichment or depletion of alloying elements to grain boundaries under irradiation. The diffusion flux of alloying species driven by the chemical potential gradient is expressed as
\begin{equation}\label{equ:fickwithmiu}
J_i = -\frac{D_i C_i}{k_B T}\,\frac{\partial \left(\mu_0 + \Delta \mu_i(x)\right)}{\partial x}
\end{equation}
where \(J_i\) is the diffusion flux of species \(i\), \(D_i\) is the diffusion coefficient, \(C_i\) is the concentration, \(k_B\) is the Boltzmann constant, \(T\) is the absolute temperature, \(\mu_0\) is the reference chemical potential, and \(\Delta \mu_i(x)\) is the radiation-induced effective chemical potential contribution.

The fluxes satisfy the mass-balance constraint
\begin{equation}\label{equ:massbalance}
\sum_i J_i + J_V = 0
\end{equation}
where \(J_V\) is the vacancy flux, given by Fick's first law:
\begin{equation}\label{equ:vacfick}
J_V = -D_V \frac{\partial C_V}{\partial x}
\end{equation}
with \(D_V\) and \(C_V\) denoting the vacancy diffusion coefficient and vacancy concentration, respectively.

For irradiated 304 stainless steel (304 SS), vacancies are continuously generated in the grain interior, whereas the grain boundary serves as a strong vacancy sink. Accordingly, the vacancy concentration profile near the grain boundary is described by
\begin{equation}\label{equ:vacconcprof}
C_V(x) = C_V^{\infty}(\mathrm{dose}) - \left(C_V^{\infty}(\mathrm{dose}) - C_V^{\mathrm{GB}}\right)\exp(-kx)
\end{equation}
where \(x\) is the distance from the grain boundary, \(C_V^{\infty}(\mathrm{dose})\) is the bulk vacancy concentration at a given irradiation dose, \(C_V^{\mathrm{GB}}\) is the vacancy concentration at the grain boundary, and \(k\) is the decay constant. In the present work, \(C_V^{\mathrm{GB}} = 0\) is assumed. Thus the vacancy flux is within the form
\begin{equation}\label{equ:vacflux}
J_V = -k D_V C_V^{\infty}(\mathrm{dose}) \exp(-kx)
\end{equation}
The mass-balance constraint (Eq. \ref{equ:massbalance}) can be rewritten as:
\begin{equation}\label{equ:massbalance2}
\sum_i J_i = k D_V C_V^{\infty}(\mathrm{dose}) \exp(-kx)
\end{equation}
By intergrating the two sides of \ref{equ:massbalance2}, the mass-balance constraint for RIS effective chemical potential can be expressed as:
\begin{equation} \label {equ:rischem}
\sum_i D_i C_i \Delta \mu_i(x) = K_B T D_V C_V^{\infty}(\mathrm{dose}) \exp(-kx) + Const
\end{equation}
We assume that the effective chemical potential is zero at infinite distance from the grain boundary, thus the constant term in Eq. \ref{equ:rischem} is zero. The effective chemical potential constraint by mass balance can be expressed as :
\begin{equation} \label{equ:risfinal}
\sum_i D_i C_i \Delta \mu_i(x) = K_B T D_V C_V^{\infty}(\mathrm{dose}) \exp(-kx)
\end{equation}
The effective chemical potential contribution for each alloying element is assumed to follow an exponential decay form:
\begin{equation}\label{equ:chemipot}
\Delta \mu_i(x) = \Delta \mu_i^{\mathrm{GB}} \exp\left(-\frac{x}{\lambda}\right)
\end{equation}
where \(\Delta \mu_i^{\mathrm{GB}}\) is the effective chemical potential contribution at the grain boundary, and \(\lambda\) is the characteristic length scale. For 304 SS,  the  segregation of Ni and Cr are mainly contributed by Inverse Kirkendall effect, \(\Delta \mu_i^{\mathrm{GB}}\) is given by:
\begin{equation}\label{equ:muCrNigb}
\Delta \mu_i^{\mathrm{GB}} = \Delta \mu_0^{\mathrm{GB}} \ln \frac{D_i}{D_{Fe}} 
\end{equation}
where \(D_{Fe}\) is the diffusion coefficient of Fe, and \(\Delta \mu_0^{\mathrm{GB}}\) is a fitting parameter that represents the overall magnitude of the effective chemical potential at the grain boundary. \(D_{Cr}>D_{Fe}>D_{Ni}\), therefore \(\Delta \mu_{Cr}^{\mathrm{GB}}>0\) and \(\Delta \mu_{Ni}^{\mathrm{GB}}<0\), which captures the depletion of Cr and enrichment of Ni at the grain boundary under irradiation. For Si, the segregation behavior is treated separately, because in addition to the vacancy-mediated inverse Kirkendall effect, Si can also migrate through the formation of solute-interstitial complexes. To account for this additional mechanism, the effective chemical potential contribution of Si is decomposed into a vacancy-related part and an interstitial-related part:
\begin{equation}\label{equ:muSiDecomp}
\Delta \mu_{Si}(x) = \Delta \mu_{Si}^{V}(x) + \Delta \mu_{Si}^{I}(x)
\end{equation}
where \(\Delta \mu_{Si}^{V}(x)\) denotes the vacancy-mediated contribution and \(\Delta \mu_{Si}^{I}(x)\) denotes the contribution associated with solute-interstitial complexes.

The vacancy-mediated contribution is assumed to follow the same functional form as that of substitutional alloying elements:
\begin{equation}\label{equ:muSiV}
\Delta \mu_{Si}^{V}(x) = \Delta \mu_{Si}^{V,\mathrm{GB}} \exp\left(-\frac{x}{\lambda_V}\right)
\end{equation}
where \(\Delta \mu_{Si}^{V,\mathrm{GB}}\) is the vacancy-mediated effective chemical potential of Si at the grain boundary and \(\lambda_V\) is the corresponding characteristic length scale.

The contribution from solute-interstitial complexes is expressed as
\begin{equation}\label{equ:muSiI}
\Delta \mu_{Si}^{I}(x) = \Delta \mu_{Si}^{I,\mathrm{GB}} \exp\left(-\frac{x}{\lambda_I}\right)
\end{equation}
where \(\Delta \mu_{Si}^{I,\mathrm{GB}}\) is the effective chemical potential contribution at the grain boundary associated with interstitial-complex transport, and \(\lambda_I\) is the characteristic decay length of this contribution.

Accordingly, the total effective chemical potential for Si becomes
\begin{equation}\label{equ:muSiTotal}
\Delta \mu_{Si}(x)
=
\Delta \mu_{Si}^{V,\mathrm{GB}} \exp\left(-\frac{x}{\lambda_V}\right)
+
\Delta \mu_{Si}^{I,\mathrm{GB}} \exp\left(-\frac{x}{\lambda_I}\right)
\end{equation}

In the present model, the interstitial-complex contribution is assumed to dominate the RIS behavior of Si, i.e.,
\begin{equation}\label{equ:SiDominant}
\left|\Delta \mu_{Si}^{I,\mathrm{GB}}\right| >
\left|\Delta \mu_{Si}^{V,\mathrm{GB}}\right|
\end{equation}
so that the net driving force reproduces the experimentally observed segregation tendency of Si at grain boundaries under irradiation. For simplicity, if \(\lambda_V\) and \(\lambda_I\) are taken to be identical, Eq. \ref{equ:muSiTotal} can be rewritten as
\begin{equation}\label{equ:muSiSimplified}
\Delta \mu_{Si}(x)
=
\left(
\Delta \mu_{Si}^{V,\mathrm{GB}} + \Delta \mu_{Si}^{I,\mathrm{GB}}
\right)
\exp\left(-\frac{x}{\lambda}\right)
\end{equation}
where the combined grain-boundary effective chemical potential term is fitted to yield the correct Si segregation behavior.
\printbibliography
\end{document}
